\documentclass[aspectratio=169, 10pt]{beamer}
\usepackage{xeCJK}
\usepackage{fontspec}
\usepackage{cite}
\usepackage{listings}
\usepackage{tikz}
\usetheme{Madrid}
\usecolortheme{beaver}

% 代码块样式设置
\lstset{
    basicstyle=\small\ttfamily,
    breaklines=true,
    frame=single,
    backgroundcolor=\color{gray!10}
}

\title[TrivialTier Demo]{TrivialTier 系统演示}
\subtitle{从策略定义到配置生成的完整流程}
\author[李梓勤, 周杰]{李梓勤、周杰}
\institute[SCU]{四川大学网络空间安全学院}
\date{2025 年 12 月 23 日}

\begin{document}

% 1. 标题页
\frame{\titlepage}

% 2. 演示概述
\begin{frame}{演示概述}
    \begin{block}{演示目标}
        展示 TrivialTier 策略引擎如何从声明式策略自动生成 WireGuard 和 nftables 配置文件
    \end{block}
    
    \begin{block}{演示场景}
        \begin{itemize}
            \item 三个节点：laptop\_li, nas\_core, scu-lab-server
            \item 一个安全策略：要求补丁等级 $\geq$ 20251101 的节点才能访问存储节点
            \item 自动生成每个节点的 WireGuard 和 nftables 配置
        \end{itemize}
    \end{block}
\end{frame}

% 3. 输入：节点定义
\begin{frame}[fragile]{输入：节点属性定义}
    \begin{block}{test\_nodes.yml}
        \lstset{language=yaml, basicstyle=\tiny\ttfamily}
        \begin{lstlisting}
nodes:
  - id: "laptop_li"
    ipv4: "10.0.0.2"
    attributes:
      system:
        os: "linux"
        patch_level: 20251201  # 满足策略
      identity:
        role: "developer"
        owner: "li-ziqin"
      network:
        zone: "trusted"

  - id: "nas_core"
    ipv4: "10.0.0.100"
    attributes:
      system:
        patch_level: 20241201  # 不满足策略
      identity:
        role: "storage"

  - id: "scu-lab-server"
    ipv4: "10.0.0.1"
    attributes:
      system:
        patch_level: 20251201  # 满足策略
      identity:
        role: "server"
        \end{lstlisting}
    \end{block}
\end{frame}

% 4. 输入：策略定义
\begin{frame}[fragile]{输入：安全策略定义}
    \begin{block}{policies.yml}
        \lstset{language=yaml, basicstyle=\tiny\ttfamily}
        \begin{lstlisting}
policies:
  - id: "pol-001"
    action: ALLOW
    subject:
      all:
        - attribute: "system.os"
          operator: "=="
          value: "linux"
        - attribute: "system.patch_level"
          operator: ">="
          value: 20251101
        - attribute: "network.zone"
          operator: "in"
          value: ["trusted", "office"]
    object:
      any:
        - attribute: "identity.role"
          operator: "=="
          value: "storage"
        - attribute: "identity.owner"
          operator: "=="
          value: "li-ziqin"
    enforcement:
      layer4:
        rules:
          - { proto: "tcp", dport: 22, action: "ACCEPT" }
          - { proto: "icmp", action: "ACCEPT" }
        \end{lstlisting}
    \end{block}
\end{frame}

% 5. 执行流程
\begin{frame}{执行流程}
    \begin{block}{命令}
        \texttt{\$ python3 main.py tests/test\_nodes.yml tests/policies.yml -o demo\_output}
    \end{block}
    
    \begin{enumerate}
        \item \textbf{步骤 1：ABAC 策略评估}
            \begin{itemize}
                \item 读取节点属性和策略定义
                \item 评估每个节点对之间的访问权限
                \item 生成访问矩阵（\texttt{input.json}）
            \end{itemize}
        
        \item \textbf{步骤 2：配置生成}
            \begin{itemize}
                \item 根据访问矩阵生成 WireGuard 配置文件（\texttt{*.conf}）
                \item 根据 L4 规则生成 nftables 策略文件（\texttt{*.nft}）
            \end{itemize}
    \end{enumerate}
\end{frame}

% 6. 访问矩阵结果
\begin{frame}[fragile]{输出：访问矩阵（input.json）}
    \begin{block}{策略评估结果}
        \lstset{language=json, basicstyle=\tiny\ttfamily}
        \begin{lstlisting}
"access_matrix": [
  {
    "source": "laptop_li",
    "destination": "nas_core",
    "connectivity": true,  // 满足策略
    "l4_rules": [
      { "proto": "tcp", "ports": [22], "action": "accept" },
      { "proto": "icmp", "action": "accept" }
    ]
  },
  {
    "source": "nas_core",
    "destination": "laptop_li",
    "connectivity": false,  // 补丁等级不足
    "l4_rules": []
  },
  {
    "source": "scu-lab-server",
    "destination": "nas_core",
    "connectivity": true,  // 满足策略
    "l4_rules": [...]
  }
]
        \end{lstlisting}
    \end{block}
    
    \begin{alertblock}{关键发现}
        \texttt{nas\_core} 由于补丁等级不足（20241201 $<$ 20251101），无法作为发起者访问其他节点
    \end{alertblock}
\end{frame}

% 7. WireGuard 配置示例
\begin{frame}[fragile]{输出：WireGuard 配置（laptop\_li.conf）}
    \begin{block}{生成的 WireGuard 配置}
        \lstset{language=bash, basicstyle=\small\ttfamily}
        \begin{lstlisting}
[Interface]
# Node: laptop_li
Address = 10.0.0.2/32
ListenPort = 51820

[Peer]
# Peer: nas_core (outbound connection)
PublicKey = 9xY2...p8Z1=
Endpoint = 5.6.7.8:51820
AllowedIPs = 10.0.0.100/32

[Peer]
# Peer: scu-lab-server (outbound connection)
PublicKey = peer_pubkey_A
Endpoint = 10.0.0.1:51820
AllowedIPs = 10.0.0.1/32
        \end{lstlisting}
    \end{block}
    
    \begin{itemize}
        \item 仅为策略允许的节点对建立 Peer 关系
        \item 自动配置端点地址和允许的 IP 范围
    \end{itemize}
\end{frame}

% 8. nftables 配置示例
\begin{frame}[fragile]{输出：nftables 配置（nas\_core.nft）}
    \begin{block}{生成的 nftables 策略}
        \lstset{language=bash, basicstyle=\tiny\ttfamily}
        \begin{lstlisting}
table inet wireguard_policy {
    chain input {
        type filter hook input priority filter; 
        policy drop;
        
        # Allow established connections
        ct state established,related accept
        
        # Allow TCP port 22 from laptop_li (10.0.0.2)
        ip saddr 10.0.0.2 tcp dport { 22 } accept
        
        # Allow ICMP from laptop_li
        ip saddr 10.0.0.2 icmp type echo-request accept
        ip saddr 10.0.0.2 icmp type echo-reply accept
        
        # Allow TCP port 22 from scu-lab-server (10.0.0.1)
        ip saddr 10.0.0.1 tcp dport { 22 } accept
        
        # Allow ICMP from scu-lab-server
        ip saddr 10.0.0.1 icmp type echo-request accept
        ip saddr 10.0.0.1 icmp type echo-reply accept
    }
}
        \end{lstlisting}
    \end{block}
    
    \begin{itemize}
        \item 默认拒绝（policy drop），仅允许策略明确的流量
        \item 基于源 IP 地址和端口号的细粒度控制
    \end{itemize}
\end{frame}

% 9. 演示总结
\begin{frame}{演示总结}
    \begin{block}{演示要点}
        \begin{itemize}
            \item \textbf{声明式策略：} 通过 YAML 定义安全策略，无需手动编写网络配置
            \item \textbf{自动化生成：} 从策略到配置的完整自动化流程
            \item \textbf{属性驱动：} 基于节点属性（如补丁等级）动态决定访问权限
            \item \textbf{双重执行：} L3（WireGuard）和 L4（nftables）协同工作
        \end{itemize}
    \end{block}
    
    \begin{block}{实际效果}
        \begin{itemize}
            \item \texttt{laptop\_li} 和 \texttt{scu-lab-server} 可以访问 \texttt{nas\_core}
            \item \texttt{nas\_core} 由于补丁等级不足，无法主动发起连接
            \item 所有配置自动生成，无需手动配置
        \end{itemize}
    \end{block}
\end{frame}

% 10. 谢谢页
\begin{frame}
    \centering
    \Huge \textcolor{blue}{谢谢！} \\
    \vspace{20pt}
    \large 欢迎提问与讨论。
\end{frame}

\end{document}

